\documentclass{article}
\usepackage{amsmath,amssymb}
\usepackage{hyperref}

% Upload this UTF-8 file through Backoffice > Knowledge > Upload document.
% This is a broad regression sample, not an exhaustive LaTeX language test.
\title{\textbf{**Equations for Lesson 3: Upload Coverage**}}
\author{}
\date{}

\begin{document}
\maketitle
\vspace{-2em}

% 1. Document structure: expect Markdown headings and paragraphs.
\section{Document structure}\label{sec:structure}
\subsection{Subsection}
\subsubsection{Subsubsection}
\paragraph{Run-in paragraph heading}First paragraph.\par Second paragraph.
\subparagraph{Small heading}Additional text.\newline A new line.

% 2. Inline formatting and escapes: expect readable Markdown text.
\section{Text formatting}
This is \textbf{bold}, \textit{italic}, \emph{emphasis},
\textsc{small caps}, \underline{underlined}, and \texttt{monospace}.
Escaped characters: 25\%, A\&B, \{braces\}, item\_one, and \#tag.
\noindent A paragraph without indentation.\quad Spaced text.
\smallskip

% 3. Links: expect Markdown links.
\section{Links}
Visit \href{https://example.com/lesson}{the lesson page} or \url{https://example.com}.

% 4. Inline and display math: expect KaTeX-renderable formulas.
\section{Formula delimiters}
Inline dollar math: $E=mc^2$ and $\vec{u}=[4,\,1]$.
Inline parenthesis math: \(a^2+b^2=c^2\).
\[
  \frac{dy}{dt}=f(x,u,t)
\]
\begin{equation}
  \int_0^1 x^2\,dx=\frac{1}{3}
\end{equation}
\begin{displaymath}
  \sum_{i=1}^{n}i=\frac{n(n+1)}{2}
\end{displaymath}

% 5. Alignment: expect aligned/gathered inside $$ and intact \\ line breaks.
\section{Multiline equations}
\begin{align*}
\vec{u} &= [4,\,1], \quad \vec{v}=[2,\,3] \\
-1\cdot\vec{v} &= [-2,\,-3] \\
\vec{u}-\vec{v} &= \vec{u}+(-1\cdot\vec{v})=[2,\,-2]
\end{align*}
\begin{gather*}
a+b=b+a \\
(a+b)+c=a+(b+c)
\end{gather*}
\begin{multline*}
(x_1+x_2+x_3+x_4) \\
= (x_1+x_2)+(x_3+x_4)
\end{multline*}

% 6. Lists: expect Markdown bullets and numbered items.
\section{Lists}
\begin{itemize}
\item First bullet.
\item Second bullet with inline math $x+y$.
  \begin{enumerate}
  \item First nested step.
  \item Second nested step.
  \end{enumerate}
\end{itemize}

% 7. Containers: expect readable text without layout wrappers.
\begin{abstract}
This abstract summarizes the formula examples.
\end{abstract}
\begin{center}
Centered source text.
\end{center}

% 8. References, footnotes, quotes and simple tables should be readable.
% Exact LaTeX numbering and footnote placement are not preserved.
\section{Review-required examples}
See Section~\ref{sec:structure}. This \footnote{Footnote content should remain visible.}
uses a footnote command.
\begin{quote}
Quoted source text should become a Markdown blockquote.
\end{quote}
\begin{tabular}{ll}
Term & Meaning \\
ODE & Ordinary differential equation
\end{tabular}

\end{document}
